\section{R ile çözümlü uygulama}
\label{sec:r-b14}

\textbf{Soru.} Aynı 24 kişinin son--ön farkları için $\bar d=-3{,}2$ ve
$s_d=6$ verilmiştir. Uygun testi seçin, yüzde 95 güven aralığını ve
standartlaştırılmış değişimi hesaplayıp sonucu raporlayın.

\begin{lstlisting}[style=prismR]
cift_sayisi <- 24
ortalama_fark <- -3.2
fark_sapmasi <- 6
standart_hata <- fark_sapmasi / sqrt(cift_sayisi)
serbestlik <- cift_sayisi - 1
t_degeri <- ortalama_fark / standart_hata
p_degeri <- 2 * pt(abs(t_degeri), df = serbestlik,
                   lower.tail = FALSE)
aralik <- ortalama_fark + c(-1, 1) *
  qt(0.975, df = serbestlik) * standart_hata
sonuc <- data.frame(
  fark = ortalama_fark, se = standart_hata,
  t = t_degeri, sd = serbestlik, p = p_degeri,
  alt = aralik[1], ust = aralik[2],
  dz = ortalama_fark / fark_sapmasi
)
print(sonuc, digits = 5)
stopifnot(p_degeri < 0.05, all(aralik < 0))
\end{lstlisting}

\begin{outputguidebox}
\textbf{Çözüm ve kontrol değerleri.} Tasarım eşleştirilmiştir; farklar
üzerinden tek örneklem $t$ hesabı yapılır. $\SE=1{,}2247$,
$t(23)=-2{,}6128$, $p=0{,}01556$, yüzde 95 güven aralığı
$[-5{,}7336;-0{,}6664]$ ve $d_z=-0{,}5333$ bulunur.

\textbf{Örnek rapor.} ``Son ölçümde puanlar ortalama 3,2 puan daha düşüktü;
$t(23)=-2{,}61$, $p=0{,}016$, yüzde 95 GA $[-5{,}73;-0{,}67]$,
$d_z=-0{,}53$.''
\end{outputguidebox}

\textbf{Yorum.} 24 bağımsız kişi çifti vardır; 48 bağımsız kişi yoktur.
Fark dağılımındaki aykırı değerler ve yaklaşık normallik ham veriden
incelenmelidir. Azalmanın olumlu mu olumsuz mu olduğu ölçeğin anlamına bağlıdır;
kontrol grubu yoksa değişim yalnız uygulamaya bağlanamaz.

\textbf{Pekiştirme ve yanıt.} Son ölçümden ön ölçüme geçildiğinde farkın
ve $t$'nin işareti değişir, çift yönlü $p$ aynı kalır. Sınav yanıtında yöntem,
işaret, belirsizlik ve bağlam birlikte verilmelidir.